%  LaTeX support: latex@mdpi.com 
%  For support, please attach all files needed for compiling as well as the log file, and specify your operating system, LaTeX version, and LaTeX editor.

%=================================================================
\documentclass[journal,article,submit,moreauthors]{Definitions/mdpi} 
%\documentclass[preprints,article,submit,moreauthors]{Definitions/mdpi} 
% For posting an early version of this manuscript as a preprint, you may use "preprints" as the journal. Changing "submit" to "accept" before posting will remove line numbers.

% Below journals will use APA reference format:
% admsci, aieduc, behavsci, businesses, econometrics, economies, education, ejihpe, famsci, games, humans, ijcs, ijfs, investment, fclam, jintelligence, journalmedia, jrfm, jsam, languages, opermanag, peacestud, pmh, psycholint, publications, tourismhosp, youth

% Below journals will use Chicago reference format:
% arts, genealogy, histories, humanities, laws, literature, religions, risks, socsci

%--------------------
% Class Options:
%--------------------
%----------
% journal
%----------
% Choose between the following MDPI journals:
% accountaudit, acn, acoustics, actuators, addictions, addictprev, adhesives, adolescents, aerobiology, aerospace, agriculture, agriengineering, agrochemicals, agronomy, ai, aichem, aieng, aimater, aimed, aipa, air, aisens, aisoc, algorithms, allergies, alloys, amh, analog, analytica, analytics, anatomia, anesthres, animals, antibiotics, antibodies, antioxidants, applbiosci, appliedchem, appliedmath, appliedphys, applmech, applmicrobiol, applnano, applsci, aps, aquacj, archaeolstud, architecture, arm, arthropoda, asc, asi, astronautics, astronomy, atmosphere, atoms, audiolres, automation, axioms, bacteria, batteries, bdcc, beverages, biochem, bioengineering, biologics, biology, biomass, biomechanics, biomed, biomedicines, biomedinformatics, biomimetics, biomolecules, biophysica, bioresourbioprod, biosensors, biosphere, biotech, birds, blockchains, bloods, blsf, brainsci, breath, buildings, cancers, carbon, cardio, cardiogenetics, cardiovascmed, catalysts, cells, ceramics, challenges, chemengineering, chemistry, chemosensors, chemproc, children, chips, cimb, civileng, cleantechnol, climate, clinbioenerg, clinpract, clockssleep, cmd, cmtr, coasts, coatings, colloids, colorants, commodities, complexities, complications, compounds, computation, computers, condensedmatter, conservation, constrmater, cosmetics, covid, crops, cryo, cryptography, crystals, csmf, ctn, culture, curroncol, cyber, dairy, data, ddc, dentistry, dermato, dermatopathology, designs, devices, dhi, diabetology, diagnostics, dietetics, digital, disabilities, diseases, diversity, dna, drones, drugreact, drylands, dynamics, earth, ebj, ecm, ecologies, edm, eesp, electricity, electrochem, electronicmat, electronics, embodiedintell, encyclopedia, endocrines, energies, eng, engproc, entomology, entropy, environments, environremediat, epidemiologia, epigenomes, esa, est, fermentation, fibers, fintech, fire, fishes, fluids, foodeng, foods, forecasting, forensicsci, forests, fossstud, foundations, fractalfract, freshwater, fuels, future, futureinternet, futurepharmacol, futurephys, futuretransp, galaxies, gases, gastroent, gastrointestdisord, gastronomy, gels, genes, geochem, geographies, geohazards, geomatics, geometry, geosciences, geotechnics, geriatrics, germs, glacies, grainsci, grasses, green, greenhealth, gucdd, gynecolcancer, hardware, hazardousmatters, healthcare, healthcmater, hearts, hemato, hematolrep, hep, heritage, higheredu, highthroughput, horticulturae, hospitals, hydrobiology, hydrogen, hydrology, hydrometeorology, hydropower, hygiene, idr, iic, ijem, ijerph, ijgi, ijmd, ijms, ijns, ijom, ijp, ijpb, ijt, ijtm, ijtpp, ijtst, ime, immuno, industries, inflammj, informatics, information, infrastructures, inorganics, insects, instruments, inventions, iot, j, jaestheticmed, jal, japma, jcdd, jcm, jcp, jcrm, jcs, jcto, jdad, jdb, jdream, jemr, jeta, jfb, jfmk, jfth, jgbg, jgg, jimaging, jlpea, jmahp, jmmp, jmms, jmp, jmse, jne, jnt, jof, joi, joitmc, joma, jop, joptm, jor, jox, jpbi, jphytomed, jpm, jsan, jtaer, jvd, jzbg, kidneydial, kinasesphosphatases, knowledge, labmed, laboratories, lam, land, lasers, legumes, life, lights, limnolrev, lipidology, liquids, livers, logics, logistics, lubricants, lymphatics, machines, macromol, magnetism, magnetochemistry, make, marinedrugs, maritime, materials, materproc, mathematics, mca, mcb, measurements, medicina, medicines, medsci, membranes, merits, metabolites, metals, meteorology, methane, metrics, metrology, micro, microarrays, microbiolres, microelectronics, micromachines, microorganisms, microplastics, microwave, minerals, mining, mmphys, modelling, molbank, molecules, mollusca, mps, msf, mti, multimedia, muscles, nanoenergyadv, nanomanufacturing, nanomaterials, ncrna, ndt, network, neuroglia, neuroimaging, neurolint, neurosci, nitrogen, notspecified, nursrep, nutraceuticals, nutrients, obesities, occuphealth, oceans, ohbm, onco, optics, oral, organics, organoids, osteology, oxygen, pandemics, parasites, parasitologia, particles, pathogens, pathophysiology, pediatrrep, pets, pharmaceuticals, pharmaceutics, pharmacoepidemiology, pharmacy, phc, philosophies, photochem, photonics, photovoltaics, phycology, physchem, physics, physiologia, plants, plasma, platforms, pollutants, polymers, polysaccharides, populations, poultry, powders, precision, precisoncol, preprints, proceedings, processes, prosthesis, proteomes, psf, psychiatryint, psychoactives, purification, quantumrep, quaternary, qubs, radiation, rareearths, rdt, reactions, realestate, receptors, recycling, regeneration, remotesensing, reports, reprodmed, resources, rheumato, rjpm, robotics, rsee, ruminants, safety, sanpp, sca, sci, scipharm, sclerosis, seeds, sensors, separations, sexes, shi, signals, sinusitis, siuj, skins, smartcities, smartfish, sna, societies, software, soilsystems, solar, solids, spectroscj, sports, standards, stats, std, stratsediment, stresses, strokerep, superintelligence, surfaces, surgeries, suschem, sustainability, sustainfoods, sustainpackag, symmetry, synbio, systems, tae, targets, taxonomy, technologies, telecom, test, textiles, thalassrep, therapeutics, thermo, timespace, tomography, toxics, toxins, tph, transplantology, transportation, traumacare, traumas, tri, tropicalmed, universe, urbansci, uro, vaccines, vehicles, venereology, vetsci, vibration, virtualworlds, viruses, vision, waste, water, welding, wem, wevj, wild, wind, women, world, zoonoticdis

%---------
% article
%---------
% The default type of manuscript is "article", but can be replaced by: 
% abstract, addendum, article, benchmark, book, bookreview, briefcommunication, briefreport, casereport, changes, clinicopathologicalchallenge, comment, commentary, communication, conceptpaper, conferenceproceedings, correction, conferencereport, creative, datadescriptor, discussion, entry, expressionofconcern, extendedabstract, editorial, essay, erratum, fieldguide, hypothesis, interestingimages, letter, meetingreport, monograph, newbookreceived, obituary, opinion, proceedingpaper, projectreport, reply, retraction, review, perspective, protocol, shortnote, studyprotocol, supfile, systematicreview, technicalnote, viewpoint, guidelines, registeredreport, tutorial,  giantsinurology, urologyaroundtheworld
% supfile = supplementary materials

%----------
% submit
%----------
% The class option "submit" will be changed to "accept" by the Editorial Office when the paper is accepted. This will only make changes to the frontpage (e.g., the logo of the journal will get visible), the headings, and the copyright information. Also, line numbering will be removed. Journal info and pagination for accepted papers will also be assigned by the Editorial Office.

%------------------
% moreauthors
%------------------
% If there is only one author the class option oneauthor should be used. Otherwise use the class option moreauthors.

%=================================================================
% MDPI internal commands - do not modify
\firstpage{1} 
\makeatletter 
\setcounter{page}{\@firstpage} 
\makeatother
\pubvolume{1}
\issuenum{1}
\articlenumber{0}
\pubyear{2026}
\copyrightyear{2026}
%\externaleditor{Firstname Lastname} % More than 1 editor, please add `` and '' before the last editor name
\datereceived{ } 
\daterevised{ } % Comment out if no revised date
\dateaccepted{ } 
\datepublished{ } 
%\datecorrected{} % For corrected papers: "Corrected: XXX" date in the original paper.
%\dateretracted{} % For retracted papers: "Retracted: XXX" date in the original paper.
%\doinum{} % Used for some special journals, like molbank
%\pdfoutput=1 % Uncommented for upload to arXiv.org
%\CorrStatement{yes}  % For updates
%\longauthorlist{yes} % For many authors that exceed the left citation part
%\IsAssociation{yes} % For association journals

%=================================================================
% Add packages and commands here. The following packages are loaded in our class file: fontenc, inputenc, calc, indentfirst, fancyhdr, graphicx, epstopdf, lastpage, ifthen, float, amsmath, amssymb, lineno, setspace, enumitem, mathpazo, booktabs, titlesec, etoolbox, tabto, xcolor, colortbl, soul, multirow, microtype, tikz, totcount, changepage, attrib, upgreek, array, tabularx, pbox, ragged2e, tocloft, marginnote, marginfix, enotez, amsthm, natbib, hyperref, cleveref, scrextend, url, geometry, newfloat, caption, draftwatermark, seqsplit
% cleveref: load \crefname definitions after \begin{document}

%=================================================================
% Please use the following mathematics environments: Theorem, Lemma, Corollary, Proposition, Characterization, Property, Problem, Example, ExamplesandDefinitions, Hypothesis, Remark, Definition, Notation, Assumption
%% For proofs, please use the proof environment (the amsthm package is loaded by the MDPI class).

%=================================================================
% Full title of the paper (Capitalized)
\Title{Arc-FaultNet: A Lightweight Dual-Branch CNN with Channel Attention and Cross-Attention for Generalizable Series Arc Fault Detection
}

% Author Orchid ID: enter ID or remove command
\newcommand{\orcidauthorA}{0000-0000-0000-000X} % Add \orcidA{} behind the author's name
%\newcommand{\orcidauthorB}{0000-0000-0000-000X} % Add \orcidB{} behind the author's name

% Authors, for the paper (add full first names)
\Author{Salim Gouaied
 $^{1}$\orcidA{}and Patrick Schweitzer $^{2}$\orcidB{} }

%\longauthorlist{yes}

% MDPI internal command: Authors, for metadata in PDF
\AuthorNames{Salim Gouaied and Patrick Schweitzer}

%\longauthorlist{yes}

% Affiliations / Addresses (Add [1] after \address if there is only one affiliation.)
\address{%
$^{1}$ \quad Institut Jean Lamour, University of Lorraine, Nancy, France, salimgouaied@gmail.com\\
$^{2}$ \quad Institut Jean Lamour, University of Lorraine, Nancy, France; patrick.schweitzer@univ-lorraine.fr}

% Contact information of the corresponding author
% \corres{Correspondence: e-mail@e-mail.com; Tel.: (optional; include country code; if there are multiple corresponding authors, add author initials) +xx-xxxx-xxx-xxxx (F.L.)}

% Current address and/or shared authorship
%\firstnote{Current address: Affiliation.}  
% Current address should not be the same as any items in the Affiliation section.

%\secondnote{These authors contributed equally to this work.}
% The commands \thirdnote{} till \eighthnote{} are available for further notes.

%\simplesumm{} % Simple summary

%\conference{} % An extended version of a conference paper

% Abstract (Do not insert blank lines, i.e. \\) 
\abstract{Series arc faults in low-voltage electrical installations pose severe fire hazards due to their intermittent, load-dependent nature, making reliable detection across diverse conditions a persistent challenge. This paper presents Arc-FaultNet, a lightweight dual-branch convolutional neural network that jointly exploits temporal and spectral representations of the line current through complementary attention mechanisms. The temporal branch extracts four physically derived channels via a 1D convolutional stack enhanced with Squeeze-and-Excitation (SE) attention, while the spectral branch processes log-power STFT spectrograms through a learnable frequency gate. Both representations are fused via a cross-conditioned channel attention mechanism enabling mutual temporal–spectral guidance. To assess generalization capacity and architectural stability, several training protocols were conducted — including strict GroupKFold cross-validation and single-model training — consistently yielding strong results: up to 98.77\% accuracy 98.68\% F1-score and 99\% under single-model training, and 90.16\% accuracy and 94.57\% specificity  under cross-validation, where cross-attention fusion outperforms naive concatenation by +5.62 pp F1. An enhanced variant with SE blocks and a deep classifier head further reduces performance variance by 28–51\%. The consistently strong results across all protocols confirm the generalization capacity and architectural robustness of Arc-FaultNet. With fewer than 365K parameters, Arc-FaultNet offers a practical path toward IEC 62606-compliant embedded deployment..
}

% Keywords
\keyword{series arc fault detection; channel attention; channel attention; squeeze-and-excitation; dual-branch CNN; STFT spectrogram;} 

% The fields PACS, MSC, and JEL may be left empty or commented out if not applicable
%\PACS{J0101}
%\MSC{}
%\JEL{}

%%%%%%%%%%%%%%%%%%%%%%%%%%%%%%%%%%%%%%%%%%
% Only for the journal Diversity
%\LSID{\url{http://}}

%%%%%%%%%%%%%%%%%%%%%%%%%%%%%%%%%%%%%%%%%%
% Only for the journal Applied Sciences
%\featuredapplication{Authors are encouraged to provide a concise description of the specific application or a potential application of the work. This section is not mandatory.}
%%%%%%%%%%%%%%%%%%%%%%%%%%%%%%%%%%%%%%%%%%

%%%%%%%%%%%%%%%%%%%%%%%%%%%%%%%%%%%%%%%%%%
% Only for the journal Data
%\dataset{DOI number or link to the deposited data set if the data set is published separately. If the data set shall be published as a supplement to this paper, this field will be filled by the journal editors. In this case, please submit the data set as a supplement.}
%\datasetlicense{License under which the data set is made available (CC0, CC-BY, CC-BY-SA, CC-BY-NC, etc.)}

%%%%%%%%%%%%%%%%%%%%%%%%%%%%%%%%%%%%%%%%%%
% Only for the journal BioTech, Fishes, Neuroimaging and Toxins
%\keycontribution{The breakthroughs or highlights of the manuscript. Authors can write one or two sentences to describe the most important part of the paper.}

%%%%%%%%%%%%%%%%%%%%%%%%%%%%%%%%%%%%%%%%%%
% Only for the journal Encyclopedia
%\encyclopediadef{For entry manuscripts only: please provide a brief overview of the entry title instead of an abstract.}

%%%%%%%%%%%%%%%%%%%%%%%%%%%%%%%%%%%%%%%%%%
% Different journals have different requirements. Please check the specific journal guidelines in the "Instructions for Authors" on the journal's official website.
%\addhighlights{yes}
%\renewcommand{\addhighlights}{%
%
%\noindent This is an optional section, the goal of which is to increase the discoverability and readability of the article via search engines and other scholars. Highlights should not be a copy of the abstract, but a simple overview allowing the reader to quickly and simply find out what the article is about and what can be cited from it. Each of these parts should be up to 2~bullet points.\vspace{3pt}\\
%\textbf{What are the main findings?}
% \begin{itemize}[labelsep=2.5mm,topsep=-3pt]
% \item First bullet.
% \item Second bullet.
% \end{itemize}\vspace{3pt}
%\textbf{What are the implications of the main findings?}
% \begin{itemize}[labelsep=2.5mm,topsep=-3pt]
% \item First bullet.
% \item Second bullet.
% \end{itemize}
%}

%%%%%%%%%%%%%%%%%%%%%%%%%%%%%%%%%%%%%%%%%%


% The order of the section titles is different for some journals. Please refer to the "Instructions for Authors” on the journal homepage.
\begin{document}

\section{Introduction}

Electrical safety in residential installations is a major industry concern for many years.
Among the faults that need to be detected, electrical arcing faults represent a significant
challenge that remains unresolved at this time. In Europe, the IEC 62606 standard defines the
tests and requirements related to the development of these devices (AFCI). An additional
challenge in developing detection systems with very high detection performance is the need to
minimize or even eliminate the false positive rate.
Many detection methods are based on the mathematical analysis of line current (transforms
associated with descriptors), with classification performed using artificial intelligence
methods. Seuls les travaux les plus récents sont cités dans la suite du document.
Yang Zhang and al. made a review of identification features and criteria for arc fault
detection. The method of detection they propose is first an event-based load classification
algorithm., They use descriptors on the current that combine frequency-domain Fast Fourier
Transform (FFT) harmonics with time-domain statistical metrics, such as Mean, Variance,
Maximum Slip Difference (MSD), and Zero Current Period (ZCP). Then, a Sequential
Floating Forward Selection (SFFS) algorithm extracts an optimal subset for each specific load
category.
Xiaofei Zhang and al. proposed a method using a periodic background subtraction technique
to eliminate normal current variations and environmental noise. They use 23 signal
descriptors across the time, frequency, and time-frequency domains to fully characterize the
current.
Liu and al. uses 8 descriptors, consisting of 5 time-domain eigenvalues (such as peak and
kurtosis indices) and 3 frequency-domain eigenvalues (such as frequency variance) to form a
comprehensive feature vector. A Fuzzy C-Means clustering algorithm initially determine
cluster centers. The clustered data is then fed into a LSTM-CNN neural network.
Dowalla and al. use a time-domain analysis of both current and voltage signals to detect series
arc faults and pinpoint the specific faulty line within an electrical network. For the detection,
16 time-domain descriptors, including metrics like Maximum Euclidean Distance, Maximum
Slip Difference, and Zero Current Period. A random forest classifier to accurately distinguish
between normal operation and arc-fault conditions.
Xin Ning and al. proposed a cost-sensitive arc fault detection method that integrates multi-
feature extraction with a Conditional Batch Normalization Convolutional Neural Network
(CBN-CNN). The method is based on time, frequency, energy, and spatial domains and
extract 6 scalar descriptors (maximum, minimum, mean, variance, power spectral entropy,
and wavelet energy entropy) alongside a 2-dimensional spatial image matrix of the current
periods.
Other approaches are based on time-frequency analysis, specifically wavelet analysis [ref] .
Other authors have proposed custom architectures. Xin Ning and al. developed Arc_EffNet, a
lightweight convolutional neural network adapted from the EffNet architecture. Instead of
extracted features, the network directly utilizes 800 raw current sampling points—captured
over a single 20 ms power cycle at 40 kHz—as its input.
Kamal Paul and Al. proposed LArcNet, a lightweight convolutional neural network designed
to detect series AC arc faults directly from 1-D time-series raw current data. The method
employs a teacher-student Knowledge Distillation (KD) technique, where a robust teacher
model transfers its learned feature representations to a simpler student model.

Huang and al. proposed an improved lightweight convolutional neural network (ShuffleNet
V2-ECA) using raw one-dimensional main current data as its descriptors, an adapted one-
dimensional convolution and an efficient channel attention mechanism.

The use of specific descriptors is particularly well-suited for simple implementation on
embedded systems. However, current methods do not always deliver optimal detection
performance. A particularly effective strategy derived from AI is based on attention
mechanisms. These mechanisms—specifically self-attention and cross-attention—have
emerged as key architectures that enable the system to focus on relatively more important
information within a sequence of fused features, thereby contributing to improved
classification accuracy.
Xin Ning and al. proposed a lightweight network composed of a convolutional neural network
(CNN) with an attention mechanism. A Short-Time Fourier Transform (STFT) map the
current signal into the frequency domain and use the model’s attention weights to identify and
retain only the most critical frequency bands.
As for them, Qu and al. proposes a series arc fault detection method based on multimodal
feature fusion using mathematical transforms including mathematical statistics, Fourier
transform, wavelet packet transform, and continuous wavelet transform. Attention
mechanisms enable the model to focus on the most meaningful data component.
Wang and al. proposed a design a Tri-Domain Dynamic Attention (TDDA) module using
time-domain, frequency-domain, and temporal derivative information associated to a CNN)
The attention module enables them to perform enhanced feature extraction.
Won Kyu Choi and al. presented a model which uses current signal descriptors extracted from
both the time domain (raw time-series data) and the frequency domain (Fast Fourier
Transform magnitudes) and a 1D CNN. A LSTM network with an attention mechanism to
focus on sequence relevance and accurately identify the specific load type.
Gao and al. presented a lightweight arc fault detection method using current signals as raw
data, where the channel attention mechanism dynamically assigns weights across feature
channels to prioritize critical information for fault identification.
In this paper, a new model Arc-FaultNet is proposed as a lightweight dual-branch
convolutional neural network capable of performing joint temporal and spectral analysis on
line current signals. This method utilizes both Channel Attention and Cross-Attention
mechanism to refine the extraction of discriminative fault features within the network.
The main contributions of this research are as follows :
- Development of a method of detection based on spatial and temporal descriptors.
- Integration of an attention mechanism ….
- …
- Performance comparison of the attention mechanism 
Citing a journal paper \citep{ref-journal}.  Now citing a book reference \citep{ref-book1,ref-book2} or other reference types \citep{ref-unpublish,ref-url}. Please use the command \citep{ref-proceeding,ref-thesis} for the following MDPI journals, which use author--date citation: Administrative Sciences, AI in Education, Arts, Behavioral Sciences, Businesses, Econometrics, Economies, Education Sciences, European Journal of Investigation in Health, Psychology and Education, Family Sciences, Future Collections, Libraries, Archives, and Museums, Games, Genealogy, Histories, Humanities, Humans, International Journal of Cognitive Sciences, International Journal of Financial Studies, Investment, Journal of Intelligence, Journalism and Media, Journal of Risk and Financial Management, Journal of Strategic and Applied Marketing, Languages, Laws, Literature, Operations Management, Peace Studies, Project Management Horizons, Psychology International, Publications, Religions, Risks, Social Sciences, Tourism and Hospitality, Youth. 

%%%%%%%%%%%%%%%%%%%%%%%%%%%%%%%%%%%%%%%%%%
\section{Materials and Methods}

The experimental platform, illustrated in Figure 1, emulates a typical residential AC power network. The configuration consists of a 230 V, 50 Hz power supply (Domestical power network), a dedicated arc fault generator and various household appliances arranged in either a combination of single loads  or masking load masking configurations. The process to generate a series arcing fault is done using the opening contacts mode between two copper electrodes (6 millimeters diameters).

\begin{figure}[htb]
    \centering
    \includegraphics[width=0.99\linewidth]{Banc de test4.pdf}
    \caption{Experimental setup.}
    \label{fig:iec}
\end{figure}



The test setup Fig. \ref{fig:bench} comprises two distinct parts: the masking loads (Loads 1 through N connected in parallel) and the arc loads (LoadN+1 to M), connected in series with the arc generator. The voltage across the arc is measured as an indicator of arc presence, and the total circuit current is also monitored for analysis. Measurements are registered using an oscilloscope (Lecroy HDO 6104) at a sampling rate of 1MHz. Line Current measurements are made using a Lecroy AP015 current probe (75MHz bandwidth).
The set of household appliances tested includes a kettle, an halogen lamp, a vacuum cleaner, a drill, 
an hair dryer, an incandescent Lightning, a PC Computer, a jigsaw and a blender. 
The first table summarizes the configurations tested without masking loads, where series arc faults occur across various combinations of interconnected loads, alongside individual loads tested separately under single-arc conditions, yielding a total of 9,386 periods for the database.

\begin{table*}[htbp]
\centering
\caption{Experimental test configurations and load combinations without masking loads.}
\label{tab:load_combinations}
\resizebox{\textwidth}{!}{%
\begin{tabular}{lccccccccccccccccccc}
\toprule
\textbf{Load / Configuration} & \textbf{1} & \textbf{2} & \textbf{3} & \textbf{4} & \textbf{5} & \textbf{6} & \textbf{7} & \textbf{8} & \textbf{9} & \textbf{10} & \textbf{11} & \textbf{12} & \textbf{13} & \textbf{14} & \textbf{15} & \textbf{16} & \textbf{17} & \textbf{18} & \textbf{19} \\
\midrule
Vaccum Cleaner & X & X & X & X & X & X & X & X & X & X & X & X & X &   &   & X &   &   &   \\
Hair Dryer     & X & X & X & X & X & X & X & X & X & X &   &   &   &   &   &   &   &   & X \\
Kettle         & X & X & X & X & X & X & X & X & X & X & X & X &   &   & X & X & X & X & X \\
Incandescent lamp & X & X & X & X & X & X & X & X & X &   &   &   &   &   &   &   &   & X &   \\
Halogen        & X &   &   &   &   &   &   &   &   &   & X & X & X & X & X & X & X &   &   \\
Drill          &   & X & X & X & X & X & X & X & X &   & X &   &   &   &   & X & X &   &   \\
Blender        &   &   &   &   & X & X & X & X & X &   &   &   &   &   & X & X &   &   &   \\
\bottomrule
\end{tabular}%
}
\end{table*}


The second table presents the configurations evaluated in the presence of masking loads, designated as M for upstream loads placed ahead of the arc fault location, and A for downstream loads positioned across the arc path (Arc loads). The total number of additional periods in the database to 9,386 periods of current with and without arc fault.



\begin{table*}[htbp]
\centering
\caption{Experimental test configurations with masking loads ($M$ for upstream, $A$ for downstream).}
\label{tab:masking_configurations}
\resizebox{\textwidth}{!}{%
\scriptsize % <-- Vous pouvez modifier ici par \footnotesize ou \scriptsize
\begin{tabular}{lcccccccc}
\toprule
\textbf{Load / Configuration} & \textbf{1} & \textbf{2} & \textbf{3} & \textbf{4} & \textbf{5} & \textbf{6} & \textbf{7} & \textbf{8} \\
\midrule
Kettle            & M & M & A & M & M & A & M & -- \\
Incandescent lamp & A & A & M & A & A & M & A & -- \\
Halogen           & A & M & M & M & M & A & M & -- \\
Vaccum Cleaner    & M & A & A & -- & -- & -- & M & A \\
Computer          & -- & M & -- & -- & -- & -- & M & A \\
Jigsaw            & -- & -- & -- & A & A & M & A & M \\
\bottomrule
\end{tabular}%
}
\end{table*}

An initial step in labeling the main dataset (training and validation datasets) must be completed. The operational state of each point from the electrical current line signature is determined by evaluating the corresponding arc voltage. Under ideal conditions, the absence of an arc is reflected by a zero-voltage baseline, while actual arcing activity produces a distinct potential rise. A predefined voltage threshold of 13 Volts is established; segments surpassing this limit are designated as faulty. Accordingly, a binary labeling scheme is applied, assigning a value of '1' to active arc occurrences and '0' to normal operating conditions.

\subsection{Dataset Composition}
\label{sec:dataset_composition}

The dataset used in this work contains a total of \textbf{10\,860 labeled current cycles}, assembled from two experimental campaigns acquired on the bench of Figure~\ref{fig:iec}. The first campaign (referred to as the \emph{July} data) contributes \textbf{9\,386 cycles} collected over three independent acquisition sessions. The second campaign (referred to as the \emph{Othmane--Salim} data) contributes \textbf{1\,474 additional cycles}, which correspond to 80\% of that campaign; the remaining 20\% (\textbf{368 cycles}) were set aside \emph{before} the dataset was assembled and were never used for training or model selection, being reserved as a fully independent hold-out set for final evaluation. Across the 10\,860 cycles, the class distribution is 5\,898 normal and 4\,962 arcing cycles, i.e. a nearly balanced binary problem. Table~\ref{tab:dataset_composition} summarizes this composition.

At the raw level, each of the 10\,860 samples originates from a triplet of CSV files recorded by the Teledyne LeCroy oscilloscope, providing three synchronized channels:
\begin{itemize}
\item \textbf{C1} --- $V_{\text{line}}$: the mains voltage;
\item \textbf{C2} --- $V_{\text{arc}}$: the voltage measured across the arc electrodes;
\item \textbf{C3} --- $I$: the line current.
\end{itemize}
Their roles in the processing pipeline are strictly separated. The mains voltage C1 is a stable, load-independent signal and is used as the anchor for cycle segmentation. The arc voltage C2 is used only as the ground-truth oracle for the labeling procedure described above, and is \emph{never} presented as an input to the detector. The line current C3 is therefore the sole signal analyzed by the model.

\begin{table}[H]
\caption{Composition of the assembled dataset.\label{tab:dataset_composition}}
\begin{tabularx}{\textwidth}{lCC}
\toprule
\textbf{Source} & \textbf{Cycles} & \textbf{Role} \\
\midrule
July data (3 acquisition sessions)      & 9\,386  & Training corpus \\
Othmane--Salim data (80\% of campaign)   & 1\,474  & Training corpus \\
\midrule
\textbf{Assembled dataset}               & \textbf{10\,860} & Train / validation / test splits \\
\midrule
Othmane--Salim data (remaining 20\%)     & 368     & Independent hold-out (final test) \\
\bottomrule
\end{tabularx}
\noindent{\footnotesize Class balance within the 10\,860 cycles: 5\,898 normal / 4\,962 arc.}
\end{table}

\subsection{Signal Conditioning}
\label{sec:signal_conditioning}

Each recording is segmented into individual 50~Hz cycles using the mains voltage C1 rather than the current C3: because C1 is sinusoidal and load-independent, its positive-going zero crossings provide consistent cycle boundaries regardless of the connected load or of the presence of an arc, whereas the current waveform is precisely the quantity distorted by the phenomenon to be detected. Consecutive zero crossings, validated by their expected spacing, delimit one complete cycle.

Every extracted cycle is then decimated from the native sampling rate of 1~MHz (20\,000 samples per cycle) to \textbf{102.4~kHz} (\textbf{2\,048 samples per cycle}) using polyphase resampling with a Kaiser-windowed FIR anti-aliasing filter. This yields a Nyquist frequency of 51.2~kHz, which preserves the discriminative arc content --- concentrated below approximately 50~kHz --- while discarding the essentially empty high-frequency band, and reduces the memory and computational footprint by an order of magnitude, a prerequisite for embedded deployment.

Finally, each current cycle is independently normalized by a per-cycle z-score, $\hat{x} = (x - \mu_x)/\sigma_x$. This enforces amplitude invariance by construction: since different loads draw currents of very different magnitudes, absolute amplitude is a load-specific, non-transferable feature, whereas the arc signature is carried by the \emph{morphology} of the waveform (discontinuities, flat shoulders, harmonic content). Normalizing per cycle removes the amplitude shortcut and forces the model to rely on load-invariant morphological cues.

# \begin{quote}
# This is an example of a quote.
# \end{quote}

%%%%%%%%%%%%%%%%%%%%%%%%%%%%%%%%%%%%%%%%%%
\section{Results}

This section may be divided by subheadings. It should provide a concise and precise description of the experimental results, their interpretation, and the experimental conclusions that can be drawn.

\subsection{Subsection}
\subsubsection{Subsubsection}

Bulleted lists look like this:
\begin{itemize}
\item	First bullet;
\item	Second bullet;
\item	Third bullet.
\end{itemize}

Numbered lists can be added as follows:
\begin{enumerate}
\item	First item; 
\item	Second item;
\item	Third item.
\end{enumerate}

The text continues here.

\subsection{Figures, Tables and Schemes}

All figures and tables should be cited in the main text as Figure~\ref{fig1}, Table~\ref{tab1}, etc.

\begin{figure}[H]
%\isPreprints{\centering}{} % Only used for preprints
\includegraphics[width=8.0 cm]{Definitions/logo-mdpi}
\caption{This is a figure. Schemes follow the same formatting.\label{fig1}}
\end{figure}   
\unskip

\begin{table}[H] 
%\small % Change table font size
\caption{This is a table caption. Tables should be placed in the main text near the first time they are~cited.\label{tab1}}
%\isPreprints{\centering}{} % Only used for preprints
\begin{tabularx}{\textwidth}{CCC}
\toprule
\textbf{Title 1}	& \textbf{Title 2}	& \textbf{Title 3}\\
\midrule
Entry 1		& Data			& Data\\
Entry 2		& Data			& Data \textsuperscript{1}\\
\bottomrule
\end{tabularx}

\noindent{\footnotesize{\textsuperscript{1} Tables may have a footer.}}
\end{table}

The text continues here (Figure~\ref{fig2} and Table~\ref{tab2}).

% Example of a figure that spans the whole page width and with subfigures. The same concept works for tables, too.
\begin{figure}[H]
%\isPreprints{}{% This command is only used for ``preprints''.
\begin{adjustwidth}{-\extralength}{0cm}
\centering
%} % If the paper is ``preprints'', please uncomment this parenthesis.
\subfloat[\centering]{\includegraphics[width=7.0cm]{Definitions/logo-mdpi}}
%\hfill
\subfloat[\centering]{\includegraphics[width=7.0cm]{Definitions/logo-mdpi}}\\
\subfloat[\centering]{\includegraphics[width=7.0cm]{Definitions/logo-mdpi}}
%\hfill
\subfloat[\centering]{\includegraphics[width=7.0cm]{Definitions/logo-mdpi}}
%\isPreprints{}{% This command is only used for ``preprints''.
\end{adjustwidth}
%} % If the paper is ``preprints'', please uncomment this parenthesis.
\caption{This is a wide figure. Schemes follow the same formatting. If there are multiple panels, they should be listed as follows: (\textbf{a}) Description of what is contained in the first panel. (\textbf{b}) Description of what is contained in the second panel. (\textbf{c}) Description of what is contained in the third panel. (\textbf{d})~Description of what is contained in the fourth panel. Figures should be placed in the main text near the first time they are cited.\label{fig2}}
\end{figure} 

\begin{table}[H]
\caption{This is a wide table. Tables should be placed in the main text near the first time they are~cited.\label{tab2}}
%\isPreprints{\centering}{% This command is only used for ``preprints''.
	\begin{adjustwidth}{-\extralength}{0cm}
%} % If the paper is ``preprints'', please uncomment this parenthesis.
%\isPreprints{\begin{tabularx}{\textwidth}{CCCC}}{% This command is only used for ``preprints''.
		\begin{tabularx}{\fulllength}{CCCC}
%} % If the paper is ``preprints'', please uncomment this parenthesis.
			\toprule
			\textbf{Title 1}	& \textbf{Title 2}	& \textbf{Title 3}     & \textbf{Title 4}\\
			\midrule
\multirow[m]{3}{*}{Entry 1 *}	& Data			& Data			& Data\\
			  	                   & Data			& Data			& Data\\
			             	      & Data			& Data			& Data\\
                   \midrule
\multirow[m]{3}{*}{Entry 2}    & Data			& Data			& Data\\
			  	                  & Data			& Data			& Data\\
			             	     & Data			& Data			& Data\\
			\bottomrule
		\end{tabularx}
%		\isPreprints{}{% This command is only used for ``preprints''.
	\end{adjustwidth}
%} % If the paper is ``preprints'', please uncomment this parenthesis.
	\noindent{\footnotesize{* Tables may have a footer.}}
\end{table}

%\begin{listing}[H]
%\caption{Title of the listing}
%\rule{\columnwidth}{1pt}
%\raggedright Text of the listing. In font size footnotesize, small, or normalsize. Preferred format: left aligned and single spaced. Preferred border format: top border line and bottom border line.
%\rule{\columnwidth}{1pt}
%\end{listing}

Text.

Text.

\subsection{Formatting of Mathematical Components}

This is example 1 of an equation:
\begin{equation}
a = 1,
\end{equation}
the text following an equation need not be a new paragraph if it continues a sentence. Please punctuate equations as regular text.

This is example 2 of an equation:
%\isPreprints{}{% This command is only used for ``preprints''.
\begin{adjustwidth}{-\extralength}{0cm}
%} % If the paper is ``preprints'', please uncomment this parenthesis.
\begin{equation}
a = b + c + d + e + f + g + h + i + j + k + l + m + n + o + p + q + r + s + t + u + v + w + x + y + z
\end{equation}
%\isPreprints{}{% This command is only used for ``preprints''.
\end{adjustwidth}
%} % If the paper is ``preprints'', please uncomment this parenthesis.

If the text following an equation starts a new sentence, it should be a new paragraph. Please punctuate equations as regular text.

%% Example of a page in landscape format (with table and table footnote).
%\startlandscape
%\begin{table}[H] %% Table in wide page
%%\isPreprints{\centering}{} % This command is only used for ``preprints''.
%\caption{This is a very wide table.\label{tab3}}
%	\begin{tabularx}{\textwidth}{CCCC}
%		\toprule
%		\textbf{Title 1}	& \textbf{Title 2}	& \textbf{Title 3}	& \textbf{Title 4}\\
%		\midrule
%		Entry 1		& Data			& Data			& This cell has some longer content that runs over two lines.\\
%		Entry 2		& Data			& Data			& Data\textsuperscript{1}\\
%		\bottomrule
%	\end{tabularx}
%%\isPreprints{}{% This command is only used for ``preprints''.
%	\begin{adjustwidth}{+\extralength}{0cm}
%%} % If the paper is ``preprints'', please uncomment this parenthesis.
%		\noindent\footnotesize{\textsuperscript{1} This is a table footnote.}
%%\isPreprints{}{% This command is only used for ``preprints''.
%	\end{adjustwidth}
%%} % If the paper is ``preprints'', please uncomment this parenthesis.
%\end{table}
%\finishlandscape

Theorem-type environments (including propositions, lemmas, corollaries, etc.) can be formatted as follows:
%% Example of a theorem:
\begin{Theorem}
Example text of a theorem. Theorems, propositions, lemmas, etc., should be numbered sequentially (i.e., Proposition 2 follows Theorem 1). Examples or Remarks use the same formatting, but should be numbered separately, so a document may contain Theorem 1, Remark 1 and Example~1.
\end{Theorem}

The text continues here. Proofs must be formatted as follows:

%% Example of a proof:
\begin{proof}[Proof of Theorem 1]
Text of the proof. Note that the phrase ``of Theorem 1” is optional if it is clear which theorem is being referred to. Always finish a proof with the following symbol.
\end{proof}

The text continues here.

%%%%%%%%%%%%%%%%%%%%%%%%%%%%%%%%%%%%%%%%%%
\section{Discussion}

Authors should discuss the results and how they can be interpreted from the perspective of previous studies and of the working hypotheses. The findings and their implications should be discussed in the broadest context possible. Future research directions may also be highlighted.

%%%%%%%%%%%%%%%%%%%%%%%%%%%%%%%%%%%%%%%%%%
\section{Conclusions}

This section is not mandatory but can be added to the manuscript if the discussion is unusually long or complex.

%%%%%%%%%%%%%%%%%%%%%%%%%%%%%%%%%%%%%%%%%%
\section{Patents}

This section is not mandatory but may be added if there are patents resulting from the work reported in this manuscript.

%%%%%%%%%%%%%%%%%%%%%%%%%%%%%%%%%%%%%%%%%%
\vspace{6pt} 

%%%%%%%%%%%%%%%%%%%%%%%%%%%%%%%%%%%%%%%%%%
%% optional
%\supplementary{The following supporting information can be downloaded at \linksupplementary{s1}, Figure S1: title; Table S1: title; Video S1: title.}

% Only for journal Methods and Protocols:
% If you wish to submit a video article, please do so with any other supplementary material.
% \supplementary{The following supporting information can be downloaded at \linksupplementary{s1}, Figure S1: title; Table S1: title; Video S1: title. A supporting video article is available at doi: link.}

% Only used for preprtints:
% \supplementary{The following supporting information can be downloaded at the website of this paper posted on \href{https://www.preprints.org/}{Preprints.org}.}

% Only for journal Hardware:
% If you wish to submit a video article, please do so with any other supplementary material.
% \supplementary{The following supporting information can be downloaded at \linksupplementary{s1}, Figure S1: title; Table S1: title; Video S1: title.\vspace{6pt}\\
%\begin{tabularx}{\textwidth}{lll}
%\toprule
%\textbf{Name} & \textbf{Type} & \textbf{Description} \\
%\midrule
%S1 & Python script (.py) & Script of python source code used in XX \\
%S2 & Text (.txt) & Script of modelling code used to make Figure X \\
%S3 & Text (.txt) & Raw data from experiment X \\
%S4 & Video (.mp4) & Video demonstrating the hardware in use \\
%... & ... & ... \\
%\bottomrule
%\end{tabularx}
%}

%%%%%%%%%%%%%%%%%%%%%%%%%%%%%%%%%%%%%%%%%%
\authorcontributions{For research articles with several authors, a short paragraph specifying their individual contributions must be provided. The following statements should be used ``Conceptualization, X.X. and Y.Y.; methodology, X.X.; software, X.X.; validation, X.X., Y.Y. and Z.Z.; formal analysis, X.X.; investigation, X.X.; resources, X.X.; data curation, X.X.; writing---original draft preparation, X.X.; writing---review and editing, X.X.; visualization, X.X.; supervision, X.X.; project administration, X.X.; funding acquisition, Y.Y. All authors have read and agreed to the published version of the manuscript.'', please turn to the  \href{http://img.mdpi.org/data/contributor-role-instruction.pdf}{CRediT taxonomy} for the term explanation. Authorship must be limited to those who have contributed substantially to the work~reported.}

\funding{Please add: ``This research received no external funding'' or ``This research was funded by NAME OF FUNDER grant number XXX''. Check carefully that the details given are accurate and use the standard spelling of funding agency names at \url{https://search.crossref.org/funding}, any errors may affect your future funding.}

\institutionalreview{In this section, please add the Institutional Review Board Statement and approval number for studies involving humans or animals. You might choose to exclude this statement if the study did not require ethical approval. Please note that the Editorial Office might ask you for further information. Please add ``The study was conducted in accordance with the Declaration of Helsinki, and approved by the Institutional Review Board (or Ethics Committee) of NAME OF INSTITUTE (protocol code XXX and date of approval).” for studies involving humans OR ``The animal study protocol was approved by the Institutional Review Board (or Ethics Committee) of NAME OF INSTITUTE (protocol code XXX and date of approval).” for studies involving animals OR ``Ethical review and approval were waived for this study due to REASON (please provide a detailed justification).” OR ``Not applicable.” for studies not involving humans or animals.}

\informedconsent{Any research article describing a study involving humans should contain this statement. Please add ``Informed consent was obtained from all subjects involved in the study.'' OR ``Patient consent was waived due to REASON (please provide a detailed justification).'' OR ``Not applicable'' for studies not involving humans. You might also choose to exclude this statement if the study did not involve humans.

Written informed consent for publication must be obtained from participating patients who can be identified (including by the patients themselves). Please state ``Written informed consent has been obtained from the patient(s) to publish this paper.” if applicable.}

\dataavailability{We encourage all authors of articles published in MDPI journals to share their research data. In this section, please provide details regarding where data supporting reported results can be found, including links to publicly archived datasets analyzed or generated during the study. Where no new data were created, or where data are unavailable due to privacy or ethical restrictions, a statement is still required. Suggested Data Availability Statements are available in the section ``MDPI Research Data Policies” at \url{https://www.mdpi.com/ethics}.} 

%\durcstatement{Current research is limited to the [please insert a specific academic field, e.g., XXX], which is beneficial [share benefits and/or primary use] and does not pose a threat to public health or national security. Authors acknowledge the dual-use potential of the research involving xxx and confirm that all necessary precautions have been taken to prevent potential misuse. As an ethical responsibility, authors strictly adhere to relevant national and international laws about DURC. Authors advocate for responsible deployment, ethical considerations, regulatory compliance, and transparent reporting to mitigate misuse risks and foster beneficial outcomes.}

% Only for journal Nursing Reports
%\publicinvolvement{Please describe how the public (patients, consumers, carers) were involved in the research. Consider reporting against the GRIPP2 (Guidance for Reporting Involvement of Patients and the Public) checklist. If the public were not involved in any aspect of the research add: ``No public involvement in any aspect of this research''.}
%
%% Only for journal Nursing Reports
%\guidelinesstandards{Please add a statement indicating which reporting guideline was used when drafting the report. For example, ``This manuscript was drafted against the XXX (the full name of reporting guidelines and citation) for XXX (type of research) research''. A complete list of reporting guidelines can be accessed via the equator network: \url{https://www.equator-network.org/}.}
%
%% Only for journal Nursing Reports
%\useofartificialintelligence{Please describe in detail any and all uses of artificial intelligence (AI) or AI-assisted tools used in the preparation of the manuscript. This may include, but is not limited to, language translation, language editing and grammar, or generating text. Alternatively, please state that “AI or AI-assisted tools were not used in drafting any aspect of this manuscript”.}

\acknowledgments{In this section, you can acknowledge any support given that is not covered by the author contribution or funding sections. This may include administrative and technical support, or donations in kind (e.g., materials used for experiments). Where GenAI has been used for purposes such as generating text, data, or graphics, or for study design, data collection, analysis, or interpretation of data, please add ``During the preparation of this manuscript/study, the author(s) used [tool name, version information] for the purposes of [description of use]. The authors have reviewed and edited the output and take full responsibility for the content of this publication.”}

\conflictsofinterest{Declare conflicts of interest or state ``The authors declare no conflicts of interest.” Authors must identify and declare any personal circumstances or interests that may be perceived as inappropriately influencing the representation or interpretation of reported research results. Any role of the funders in the design of the study; in the collection, analyses or interpretation of data; in the writing of the manuscript; or in the decision to publish the results must be declared in this section. If there is no role, please state ``The funders had no role in the design of the study; in the collection, analyses, or interpretation of data; in the writing of the manuscript; or in the decision to publish the~results”.} 

%%%%%%%%%%%%%%%%%%%%%%%%%%%%%%%%%%%%%%%%%%
%% Optional

%% Only for journal Encyclopedia
%\entrylink{The Link to this entry published on the encyclopedia platform.}

\abbreviations{Abbreviations}{%
The following abbreviations are used in this manuscript:\\

\noindent 
\begin{tabular}{@{}ll}
MDPI & Multidisciplinary Digital Publishing Institute\\
DOAJ & Directory of open access journals\\
TLA   & Three-letter acronym\\
LD      & Linear dichroism
\end{tabular}
}

%%%%%%%%%%%%%%%%%%%%%%%%%%%%%%%%%%%%%%%%%%
%% Optional
\appendixtitles{no} % Leave argument "no" if all appendix headings stay EMPTY (then no dot is printed after "Appendix A"). If the appendix sections contain a heading then change the argument to "yes".
\appendixstart
\appendix
\section[\appendixname~\thesection]{}
\subsection[\appendixname~\thesubsection]{}

The appendix is an optional section that can contain details and data supplemental to the main text---for example, explanations of experimental details that would disrupt the flow of the main text but nonetheless remain crucial to understanding and reproducing the research shown; figures of replicates for experiments of which representative data are shown in the main text can be added here if brief, or as Supplementary Materials. Mathematical proofs of results not central to the paper can be added as an appendix.

\begin{table}[H] 
\caption{This is a table caption.\label{tab5}}
%\newcolumntype{C}{>{\centering\arraybackslash}X}
\begin{tabularx}{\textwidth}{CCC}
\toprule
\textbf{Title 1}	& \textbf{Title 2}	& \textbf{Title 3}\\
\midrule
Entry 1		& Data			& Data\\
Entry 2		& Data			& Data\\
\bottomrule
\end{tabularx}
\end{table}

\section[\appendixname~\thesection]{}

All appendix sections must be cited in the main text. In the appendices, figures, tables, etc., should be labeled starting with “A”, e.g., Figure A1, Figure A2, etc.

%%%%%%%%%%%%%%%%%%%%%%%%%%%%%%%%%%%%%%%%%%
%\isPreprints{}{% This command is only used for ``preprints''.
\begin{adjustwidth}{-\extralength}{0cm}
%} % If the paper is ``preprints'', please uncomment this parenthesis.
%\printendnotes[custom] % Un-comment to print a list of endnotes

\reftitle{References}

% Please provide either the correct journal abbreviation (e.g. according to the “The list of Title Word Abbreviations” https://portal.issn.org/ltwa) or the full name of the journal. 
% Citations and references in the Supplementary Materials are permitted provided that they also appear in the reference list here. 

%=====================================
% References, variant A: external bibliography
%=====================================
% \bibliography{your_external_BibTeX_file}

%=====================================
% References, variant B: internal bibliography
%=====================================

% ACS format
\isAPAandChicago{}{%
\begin{thebibliography}{999}
% Reference 1
\bibitem[Author1(year)]{ref-journal}
Author~1, T. The title of the cited article. {\em Journal Abbreviation} {\bf 2008}, {\em 10}, 142--149.
% Reference 2
\bibitem[Author2(year)]{ref-book1}
Author~1, L. The title of the cited contribution. In {\em The Book Title}; Editor 1, F., Editor 2, A., Eds.; Publishing House: City, Country, 2007; pp. 32--58.
% Reference 3
\bibitem[Author3 and Author4 (year)]{ref-book2}
Author 1, A.; Author 2, B. \textit{Book Title}, 3rd ed.; Publisher: Publisher Location, Country, 2008; pp. 154--196.
% Reference 4
\bibitem[Author5 and Author6(year)]{ref-unpublish}
Author 1, A.B.; Author 2, C. Title of Unpublished Work. \textit{Abbreviated Journal Name} year, \textit{phrase indicating stage of publication (submitted; accepted; in press)}.
% Reference 5
\bibitem[Author7(year)]{ref-url}
Title of Site. Available online: URL (accessed on Day Month Year).
% Reference 6
\bibitem[Author8 et al.(year)]{ref-proceeding}
Author 1, A.B.; Author 2, C.D.; Author 3, E.F. Title of presentation. In Proceedings of the Name of the Conference, Location of Conference, Country, Date of Conference (Day Month Year); Abstract Number (optional), Pagination (optional).
% Reference 7
\bibitem[Author9(year)]{ref-thesis}
Author 1, A.B. Title of Thesis. Level of Thesis, Degree-Granting University, Location of University, Date of Completion.
\end{thebibliography}
}

% Chicago format (Used for journal: arts, genealogy, histories, humanities, laws, literature, religions, risks, socsci)
\isChicagoStyle{%
\begin{thebibliography}{999}
% Reference 1
\bibitem[Aranceta-Bartrina(1999a)]{ref-journal}
Aranceta-Bartrina, Javier. 1999a. Title of the cited article. \textit{Journal Title} 6: 100--10.
% Reference 2
\bibitem[Aranceta-Bartrina(1999b)]{ref-book1}
Aranceta-Bartrina, Javier. 1999b. Title of the chapter. In \textit{Book Title}, 2nd ed. Edited by Editor 1 and Editor 2. Publication place: Publisher, vol. 3, pp. 54–96.
% Reference 3
\bibitem[Baranwal and Munteanu {[1921]}(1955)]{ref-book2}
Baranwal, Ajay K., and Costea Munteanu. 1955. \textit{Book Title}. Publication place: Publisher, pp. 154--96. First published 1921 (op-tional).
% Reference 4
\bibitem[Berry and Smith(1999)]{ref-thesis}
Berry, Evan, and Amy M. Smith. 1999. Title of Thesis. Level of Thesis, Degree-Granting University, City, Country. Identifi-cation information (if available).
% Reference 5
\bibitem[Cojocaru et al.(1999)]{ref-unpublish}
Cojocaru, Ludmila, Dragos Constatin Sanda, and Eun Kyeong Yun. 1999. Title of Unpublished Work. \textit{Journal Title}, phrase indicating stage of publication.
% Reference 6
\bibitem[Driver et al.(2000)]{ref-proceeding}
Driver, John P., Steffen Rohrs, and Sean Meighoo. 2000. Title of Presentation. In \textit{Title of the Collected Work} (if available). Paper presented at Name of the Conference, Location of Conference, Date of Conference.
% Reference 7
\bibitem[Harwood(2008)]{ref-url}
Harwood, John. 2008. Title of the cited article. Available online: URL (accessed on Day Month Year).
\end{thebibliography}
}{}

% APA format (Used for journal: admsci, aieduc, behavsci, businesses, econometrics, economies, education, ejihpe, famsci, games, humans, ijcs, ijfs, investment, fclam, jintelligence, journalmedia, jrfm, jsam, languages, opermanag, peacestud, pmh, psycholint, publications, tourismhosp, youth)
\isAPAStyle{%
\begin{thebibliography}{999}
% Reference 1
\bibitem[\protect\citeauthoryear{Azikiwe \BBA\ Bello}{{2020a}}]{ref-journal}
Azikiwe, H., \& Bello, A. (2020a). Title of the cited article. \textit{Journal Title}, \textit{Volume}(Issue), 
Firstpage--Lastpage/Article Number.
% Reference 2
\bibitem[\protect\citeauthoryear{Azikiwe \BBA\ Bello}{{2020b}}]{ref-book1}
Azikiwe, H., \& Bello, A. (2020b). \textit{Book title}. Publisher Name.
% Reference 3
\bibitem[Davison(1623/2019)]{ref-book2}
Davison, T. E. (2019). Title of the book chapter. In A. A. Editor (Ed.), \textit{Title of the book: Subtitle} 
(pp. Firstpage--Lastpage). Publisher Name. (Original work published 1623) (Optional).
% Reference 4
\bibitem[Fistek et al.(2017)]{ref-proceeding}
Fistek, A., Jester, E., \& Sonnenberg, K. (2017, Month Day). Title of contribution [Type of contribution]. Conference Name, Conference City, Conference Country.
% Reference 5
\bibitem[Hutcheson(2012)]{ref-thesis}
Hutcheson, V. H. (2012). \textit{Title of the thesis} [XX Thesis, Name of Institution Awarding the Degree].
% Reference 6
\bibitem[Lippincott \& Poindexter(2019)]{ref-unpublish}
Lippincott, T., \& Poindexter, E. K. (2019). \textit{Title of the unpublished manuscript} [Unpublished manuscript/Manuscript in prepara-tion/Manuscript submitted for publication]. Department Name, Institution Name.
% Reference 7
\bibitem[Harwood(2008)]{ref-url}
Harwood, J. (2008). \textit{Title of the cited article}. Available online: URL (accessed on Day Month Year).
\end{thebibliography}
}{}

% If authors have biography, please use the format below
%\section*{Short Biography of Authors}
%\bio
%{\raisebox{-0.35cm}{\includegraphics[width=3.5cm,height=5.3cm,clip,keepaspectratio]{Definitions/author1.pdf}}}
%{\textbf{Firstname Lastname} Biography of first author}
%
%\bio
%{\raisebox{-0.35cm}{\includegraphics[width=3.5cm,height=5.3cm,clip,keepaspectratio]{Definitions/author2.jpg}}}
%{\textbf{Firstname Lastname} Biography of second author}

% For the MDPI journals use author-date citation, please follow the formatting guidelines on http://www.mdpi.com/authors/references
% To cite two works by the same author: \citeauthor{ref-journal-1a} (\citeyear{ref-journal-1a}, \citeyear{ref-journal-1b}). This produces: Whittaker (1967, 1975)
% To cite two works by the same author with specific pages: \citeauthor{ref-journal-3a} (\citeyear{ref-journal-3a}, p. 328; \citeyear{ref-journal-3b}, p.475). This produces: Wong (1999, p. 328; 2000, p. 475)

%%%%%%%%%%%%%%%%%%%%%%%%%%%%%%%%%%%%%%%%%%
%% for journal Sci
%\reviewreports{\\
%Reviewer 1 comments and authors’ response\\
%Reviewer 2 comments and authors’ response\\
%Reviewer 3 comments and authors’ response
%}
%%%%%%%%%%%%%%%%%%%%%%%%%%%%%%%%%%%%%%%%%%
\PublishersNote{}
%\isPreprints{}{% This command is only used for ``preprints''.
\end{adjustwidth}
%} % If the paper is ``preprints'', please uncomment this parenthesis.
\end{document}

